\SGASection{5}{Conjugacy theorem}
\SGARefTarget{sga3:xxvi:section:5}

\begin{sgastatement}{Theorem 5.1}
Let \(S\) be a semilocal scheme, \(G\) an \(S\)-reductive group, and
\(P,P'\) two opposite parabolic subgroups of \(G\) \textup{(\SGARefLink{sga3:xxvi:definition:4-3-3}{4.3.3})}.
Then
\[
 \rad^u(P)(S)\cdot P'\cdot P=G.
\]
In other words, as \(u\) runs through \(\rad^u(P)(S)\), the open
subsets \(uP'\cdot P\) of \(G\) \textup{(\SGARefLink{sga3:xxvi:remark:4:3:6}{4.3.6})} cover all of \(G\).
\end{sgastatement}
\SGARefTarget{sga3:xxvi:theorem:5:1}

The proof is divided into several steps.

\SGAInlineRefTarget{sga3:xxvi:paragraph:5-1-1}\paragraph{\SGARefLink{sga3:xxvi:paragraph:5-1-1}{5.1.1}.}
\SGARefTarget{sga3:xxvi:item:5:1:1}
It is enough to prove the theorem when \(S\) is the spectrum of a
field \(k\); this follows at once from \SGARefLink{sga3:xxvi:cor:2:6}{2.6}.

\SGAInlineRefTarget{sga3:xxvi:paragraph:5-1-2}\paragraph{\SGARefLink{sga3:xxvi:paragraph:5-1-2}{5.1.2}.}
\SGARefTarget{sga3:xxvi:item:5:1:2}
Put \(L=P\cap P'\).  Suppose that \(L\) has a Borel subgroup \(B_L\),
and let \(T\) be a maximal torus of \(B_L\), as in Exp.\ XXII, \SGARefLink{sga3:xxii:cor:5.9.7}{5.9.7}.
One checks immediately that
\[
 B=B_L\cdot\rad^u(P)
\]
is a Borel subgroup of \(P\).  Let \(B'\) be the Borel subgroup of
\(G\) opposite to \(B\) relative to \(T\), so that \(B\cap B'=T\).
After splitting \(G\) relative to \(T\), one immediately sees that
\(B'\subset P'\).  We claim that
\[
 B^u(S)\cdot B'\cdot B
 \subset \rad^u(P)(S)\cdot P'\cdot P.                 \tag{x}\label{sga3:xxvi:theorem:5-1:equation:x}
\]
Indeed,
\[
 B^u(S)\subset P(S)
 =\rad^u(P)(S)\cdot L(S)
 \subset\rad^u(P)(S)\cdot P'(S),
\]
so it is enough to observe that \(B'\cdot B\subset P'\cdot P\).
Thus \textup{(x)} reduces \SGARefLink{sga3:xxvi:theorem:5:1}{Theorem 5.1} to the pair \((B,B')\).

\SGAInlineRefTarget{sga3:xxvi:paragraph:5-1-3}\paragraph{\SGARefLink{sga3:xxvi:paragraph:5-1-3}{5.1.3}.}
\SGARefTarget{sga3:xxvi:item:5:1:3}
The theorem holds when \(k\) is algebraically closed: the hypothesis
of \SGARefLink{sga3:xxvi:paragraph:5-1-2}{5.1.2} is then satisfied, and Exp.\ XXII, \SGARefLink{sga3:xxii:prop:5.7.10}{5.7.10} applies.

\SGAInlineRefTarget{sga3:xxvi:paragraph:5-1-4}\paragraph{\SGARefLink{sga3:xxvi:paragraph:5-1-4}{5.1.4}.}
\SGARefTarget{sga3:xxvi:item:5:1:4}
The theorem also holds when \(k\) is infinite.  Indeed, by \SGARefLink{sga3:xxvi:cor:2:7}{2.7},
\[
 \rad^u(P)(k)\ \text{is dense in}\ \rad^u(P)(\bar k),
\]
and the theorem holds over \(\bar k\).

\paragraph{5.1.5.}
\SGARefTarget{sga3:xxvi:item:5:1:5}
We are therefore reduced to the case of a finite field \(k\).
The scheme \(\Bor(L)\) is a smooth homogeneous space under \(L\);
Lang's theorem \textup{(Amer.\ J. Math.\ \textbf{78} \textup{(1956)})}
therefore shows that \(L\) has a Borel subgroup \(B_L\).
\footnote{\emph{Editor's note.} See also \textup{[DG70]},
\(\S\)III.5, \SGARefLink{sga3:xxvi:theorem:7-4}{7.4}.}
By 5.1.2, we may consequently suppose that \(P=B\) and \(P'=B'\)
are Borel subgroups.  We write \(T=B\cap B'\).

\SGAInlineRefTarget{sga3:xxvi:paragraph:5-1-6}\paragraph{5.1.6.}
\SGARefTarget{sga3:xxvi:item:5:1:6}
Let \(K\) be an algebraic closure of \(k\).  Choose a pinning of
\((G_K,B_K,T_K)\), and let \(R^+\), respectively \(\Delta\), be the
set of positive, respectively simple, roots.  By Exp.\ XXII, \SGARefLink{sga3:xxii:def:5.7.2}{5.7.2},
it is enough to prove, for every \(\alpha\in\Delta\), that
\[
 u_\alpha B'^{\,u}(K)
 \subset B^u(k)B'^{\,u}(K)B(K).                       \tag{1}\label{sga3:xxvi:theorem:5-1:equation:1}
\]
Let \(\alpha_i\) be the distinct conjugates of \(\alpha\) over \(k\).
They belong to \(\Delta\), since \(B\) is defined over \(k\).  Let
\(R'\) be the set of roots which are linear combinations of the
\(\alpha_i\), and put \(R'_- = R_-\cap R'\).  Since \(R'\) is defined
over \(k\), there is a subtorus \(Q\) of \(T\) such that \(Q_K\) is
the maximal torus of the common kernel of the \(\alpha_i\).

Put
\[
 Z=\Centr_G(Q),\qquad B_Z=B\cap Z,\qquad B'_Z=B'\cap Z
\]
\textup{(Exp.\ XXII, \SGARefLink{sga3:xxii:cor:5.10.2}{5.10.2})}.  It is enough to verify the desired
assertion in \(Z\), namely
\[
 u_\alpha B_Z'^{\,u}(K)
 \subset B_Z^u(k)B_Z'^{\,u}(K)B_Z(K).                 \tag{2}\label{sga3:xxvi:theorem:5-1:equation:2}
\]
Indeed,
\[
 (B_Z'^{\,u})_K=\prod_{\beta\in R'_-}U_\beta.
\]
Let \(R''=R-R'\), and set
\[
 V=\prod_{\beta\in R''\cap R_-}U_\beta.
\]
Then
\[
 (B'^{\,u})_K=(B_Z'^{\,u})_K\cdot V,
\]
and \((B_Z)_K\) normalizes \(V\), by Exp.\ XXII, \SGARefLink{sga3:xxii:cor:5.6.7}{5.6.7}.  From
\textup{\SGARefLink{sga3:xxvi:theorem:5-1:equation:2}{(2)}} we therefore obtain successively
\[
\begin{aligned}
 u_\alpha B'^{\,u}(K)
 &=u_\alpha B_Z'^{\,u}(K)V(K)\\
 &\subset B_Z^u(k)B_Z'^{\,u}(K)B_Z(K)V(K)\\
 &\subset B_Z^u(k)B_Z'^{\,u}(K)V(K)B_Z(K),
\end{aligned}
\]
which immediately gives \textup{\SGARefLink{sga3:xxvi:theorem:5-1:equation:1}{(1)}}.  We are thus reduced to the case
\(G=Z\), that is, to the case in which
\(\operatorname{Gal}(K/k)\) acts transitively on the simple roots.

\SGAInlineRefTarget{sga3:xxvi:paragraph:5-1-7}\paragraph{\SGARefLink{sga3:xxvi:paragraph:5-1-7}{5.1.7}.}
\SGARefTarget{sga3:xxvi:item:5:1:7}
The assertion to be proved says equivalently that \(G/B\) is the
union of the translates, under \(B^u(k)\), of the open image of
\(B'^{\,u}\).  This assertion is unchanged when \(G\) is replaced by
its adjoint group, or indeed by any group having the same adjoint
group.  We may therefore suppose that \(G\) is adjoint.

\SGAInlineRefTarget{sga3:xxvi:paragraph:5-1-8}\paragraph{\SGARefLink{sga3:xxvi:paragraph:5-1-8}{5.1.8}.}
\SGARefTarget{sga3:xxvi:item:5:1:8}
Consider the Dynkin diagram of \(G_K\).  The Galois group acts
transitively on it.  But this Galois group has only cyclic quotients,
whereas the Dynkin diagram has no cycles.  The diagram is consequently
of type \(nA_1\), with \(n\geq0\), or \(mA_2\), with \(m\geq0\).
Using the canonical decomposition of Exp.\ XXIV, \SGARefLink{sga3:xxiv:proposition:5-9}{5.9}, we may write
\[
 G=\prod_{D/k}G_0,
\]
where \(D\) is a finite \(k\)-scheme and \(G_0\) is either a torus or
is of type \(A_1\) or \(A_2\).

By Exp.\ XXIV, \SGARefLink{sga3:xxiv:corollary:5-12}{5.12}, the group \(B\) comes from a Borel subgroup
\(B_0\) of \(G_0\), and \(T\) comes from a maximal torus \(T_0\) of
\(G_0\); the group \(B'\) comes from the Borel subgroup \(B'_0\) of
\(G_0\) opposite to \(B_0\) relative to \(T_0\).  We have
\[
 B^{u}(k)=B_0^u(D),\qquad
 B'^{\,u}\cdot T\cdot B^u
 =\prod_{D/k}B_0'^{\,u}\cdot T_0\cdot B_0^u.
\]
It is therefore enough to prove the assertion for the triple
\((G_0,B_0,T_0)\).

\SGAInlineRefTarget{sga3:xxvi:paragraph:5-1-9}\paragraph{\SGARefLink{sga3:xxvi:paragraph:5-1-9}{5.1.9}.}
\SGARefTarget{sga3:xxvi:item:5:1:9}
We may consequently suppose that \(G\) has type
\(\varnothing\), \(A_1\), or \(A_2\).  Since \(G\) has a Borel
subgroup \(B\), it is quasi-split relative to \(B\), by Exp.\ XXIV,
\SGARefLink{sga3:xxiv:3.9.1}{3.9.1}, and hence split if it has type \(\varnothing\) or \(A_1\).
The theorem is already known in the split case, by Exp.\ XXII,
\SGARefLink{sga3:xxii:prop:5.7.10}{5.7.10}, so only type \(A_2\) remains.

By Exp.\ XXIV, \SGARefLink{sga3:xxiv:3.11}{3.11}, there is a morphism
\(E\to\operatorname{Spec}(k)\), a principal Galois bundle under the
group \(\mathbb Z/2\mathbb Z\) of automorphisms of the Dynkin diagram
of type \(A_2\), such that
\[
 G=\operatorname{Q\text{-}\acute{E}p}_{E/\operatorname{Spec}(k)}
 (A_2).
\]
If \(E\) has a section, \(G\) is split and the theorem follows.
Otherwise \(E=\operatorname{Spec}(k')\), where \(k'/k\) is a quadratic
extension.  Finally, as observed in \SGARefLink{sga3:xxvi:paragraph:5-1-7}{5.1.7}, we may suppose that \(G\)
is simply connected, that is, that \(G\) is a form of
\(\operatorname{SL}_{3,k}\).

\SGAInlineRefTarget{sga3:xxvi:paragraph:5-1-10}\paragraph{\SGARefLink{sga3:xxvi:paragraph:5-1-10}{5.1.10}.}
\SGARefTarget{sga3:xxvi:item:5:1:10}
We are now in the following situation: \(k\) is a finite field and
\(k'/k\) is a quadratic extension.  Pin
\(\operatorname{SL}_{3,k'}\), the group of \(3\times3\) matrices of
determinant \(1\), as follows.  Its maximal torus is the group of
diagonal matrices, its Borel subgroup is the group of upper triangular
matrices, and its root vectors are
\[
 u_\alpha=
 \begin{pmatrix}1&1&0\\0&1&0\\0&0&1\end{pmatrix},
 \qquad
 u_\beta=
 \begin{pmatrix}1&0&0\\0&1&1\\0&0&1\end{pmatrix}.
\]
The big cell \(\Omega\) is given by
\[
 \begin{pmatrix}a&b&c\\d&e&f\\g&h&i\end{pmatrix}
 \in\Omega(S)
 \quad\Longleftrightarrow\quad
 a\ \text{and}\ ae-bd\ \text{are invertible},
\]
and
\[
 B^u(k)=
 \left\{
 \left.
 \begin{pmatrix}1&\bar x&z\\0&1&x\\0&0&1\end{pmatrix}
 \ \right|\
 x,z\in k',\ z+\bar z=x\bar x
 \right\}.
\]

We must prove inclusion \textup{\SGARefLink{sga3:xxvi:theorem:5-1:equation:1}{(1)}} of \SGARefLink{sga3:xxvi:paragraph:5-1-6}{5.1.6}.  Thus, for every
\(a,b,c\in K\), we must find \(x,z\in k'\) such that
\(z+\bar z=x\bar x\) and
\[
 \begin{pmatrix}1&\bar x&z\\0&1&x\\0&0&1\end{pmatrix}
 \begin{pmatrix}1&1&0\\0&1&0\\0&0&1\end{pmatrix}
 \begin{pmatrix}1&0&0\\a&1&0\\b&c&1\end{pmatrix}
 \in\Omega(K).                                         \tag{3}\label{sga3:xxvi:theorem:5-1:equation:3}
\]
If \(a\ne-1\), take \(x=z=0\).  If \(a=-1\), the required conditions
are
\[
 \begin{cases}
 z+\bar z=x\bar x,\\
 bz-\bar x\ne0,\\
 (b+c)\bar z+bx-1\ne0.
 \end{cases}                                           \tag{4}\label{sga3:xxvi:theorem:5-1:equation:4}
\]
Let \(q\geq2\) be the number of elements of \(k\).  For every
\(m\in k\), the equation \(z+\bar z=m\), with \(z\in k'\), has
\(q\) solutions.

If \(b=0\), take \(x=1\).  It remains to solve
\[
 z+\bar z=1,\qquad c\bar z\ne1,
\]
which is always possible by the preceding observation.

If \(b\ne0\), take \(x=0\).  It remains to solve
\[
 z+\bar z=0,\qquad z\ne0,\qquad (b+c)z\ne1.
\]
This is always possible when \(q\geq3\).  When \(k=\mathbb F_2\), it
is possible if \(b+c\ne1\), by taking \(z=1\).

It remains to consider \(k=\mathbb F_2\), \(b+c=1\), and \(b\ne0\).
The system becomes
\[
 z+\bar z=x\bar x,\qquad bz\ne\bar x,\qquad
 \bar z+bx\ne1.
\]
If \(b=1\), or respectively \(b\notin k'\), take \(x=1\); the last two
conditions become
\[
 z\ne b^{-1},\qquad z\ne1-b,
\]
and follow from any solution of \(z+\bar z=1\).  Finally, if
\(b\in k'-k\), take \(x=b\) and \(z=\bar b\).  This completes the
proof.
\SGAQED
